\documentclass{article}
\usepackage[utf8]{inputenc}
\usepackage[T1]{fontenc}
\usepackage{geometry}
\usepackage{graphicx}
\usepackage{float}
\usepackage{hyperref}
\usepackage[english,indonesian]{babel}

\title{Optimasi Pemrosesan Template LaTeX Menggunakan Pendekatan Berbasis AI}
\author{Academic Support Agent}
\date{\today}

\begin{document}
\maketitle

\begin{abstract}

\end{abstract}

\chapter{PENDAHULUAN}


\chapter{TINJAUAN PUSTAKA}


% Original content (commented out):
% %Exam Template for CNY
%created by: Daniel Gómez
\documentclass[letterpaper, 11pt]{article}
\input{packages.tex}
\newcommand\myauthor{Daniel Gómez}%author
\newcommand\mycoauthor{}%coauthor
\newcommand\myinstitute{New York School}%school
\newcommand\mydepartment{Mathematics Area}%the department
\newcommand\thisterm{2}%term
\newcommand\thisevaluation{Evaluation}%name of the evaluation
\newcommand\thisgrade{9}%grade level
\newcommand\copies{120}%number of copies
\usepackage{CNYExamStyle}

\begin{document}
%first page
\thispagestyle{firststyle}
\insertidentifiers
\begin{center}
    \textbf{Criteria Assessed and final grade}\\
    \criteria[3]{Criteria 1}{Criteria 2}{Criteria 3}{Criteria 4}%Inside the brackets place the number of criteria
    %If only one criteria assessed, then the bracket may be omitted
    %inside the curly brackets place the criteria in the order you want them to appear
\end{center}
\input{Instructions.tex}%Include instructions is applicable, if not pace the percent sign before \input{Instructions.tex}, Customize your instructions as needed

%body of document   
%example
%%\begin{question}[abox true(optional if you need answerbox)]{total marks}{no. of subpoints}{box length}
%%%%\qmarks{subpoint no. of marks} insert question here
%%\end{question}
\include{Questions}
\bibliographystyle{apacite}
\bibliography{Bibliography}
\end{document}


\end{document}